% Seminar 4 - Pagina 4. Fragment pentru
% % Seminar 4 - Pagina 4. Fragment pentru
% % Seminar 4 - Pagina 4. Fragment pentru
% % Seminar 4 - Pagina 4. Fragment pentru
% \input{seminars/pages/seminar4-pagina4-ro}; fără preambul sau pagină nouă.
% NOTE EDITORIALE:
% - Enunțul introductiv pare să compare monoizii multiplicativi
%   mN+1 și tN+1 cu U(Z_m), U(Z_t), probabil grupurile elementelor
%   inversabile ale inelelor de clase de resturi. Simbolul manuscris
%   seamănă cu U mare; această interpretare rămâne provizorie.
%   Pentru enunțul inițial nu există demonstrație pe această pagină.
% - În exercițiul 9, ENUNȚUL cere element neutru la stânga ex=x
%   și invers la DREAPTA xx'=e pentru fiecare x; combinația nu este
%   suficientă pentru un grup (există contraexemplu de semigrup drept-zero).
%   DEMONSTRAȚIA manuscrisă alege însă inverse la STÂNGA x'x=e,
%   x''x'=e și deduce xx'=e și xe=x. Păstrăm separat cele două
%   lecturi, fără corectarea tacită a enunțului.
% - La exercițiul 10 trebuie adăugată ipoteza H nevidă pentru
%   caracterizarea subgrupului; demonstrația presupune h în H.
% - Manuscrisul notează elementul neutru atât cu 1, cât și cu e.
% - Nu am reconstruit pasajele barate ori urmele de pe verso.

\subsection*{Seminar 4 - Pagina 4}

\noindent\underline{Arătați că:}
\[
\begin{gathered}
 (m\mathbb N+1,\cdot)\cong(t\mathbb N+1,\cdot),
 \qquad m,t\in\mathbb N \\[2pt]
 \Longleftrightarrow\\[-2pt]
 (\mathcal U(\mathbb Z_m),\cdot)
   \cong (\mathcal U(\mathbb Z_t),\cdot).
\end{gathered}
\]
\noindent (În manuscris, $\mathcal U(\mathbb Z_m)$ este indicat drept grup.)

\medskip
\noindent\underline{Exercițiul 9.}
Fie $(G,\cdot)$ semigrup. Atunci $G$ este grup dacă și numai dacă
\begin{enumerate}
 \item există $e\in G$ astfel încât $ex=x$, pentru orice $x\in G$;
 \item pentru orice $x\in G$ există $x'\in G$ cu
       $xx'=e$ (așa cum apare în enunț).
\end{enumerate}

\noindent\underline{Demonstrație (reciproca, după manuscris).}
Fie $x\in G$. Există $x'\in G$ astfel încât $x'x=e$ și
există $x''\in G$ astfel încât $x''x'=e$. Atunci
\[
\begin{aligned}
 xx'&=(ex)x'=(x''x')xx'
       =x''(x'x)x'=x''ex'=x''x'=e,\\
 xe&=x(x'x)=(xx')x=ex=x.
\end{aligned}
\]

\medskip
\noindent\underline{Exercițiul 10.}
Fie $(G,\cdot)$ grup, $H\subseteq G$, $H$ finită. Atunci
\[
 H\leq G\quad\Longleftrightarrow\quad
 H\text{ este parte stabilă la înmulțire}.
\]

\noindent\underline{Demonstrație.}
Implicația directă este evidentă. Reciproc, fie $h\in H$ și
\[
 hH=\{hx\mid x\in H\},\qquad
 H^2=\{xy\mid x,y\in H\}.
\]
Din stabilitatea lui $H$ rezultă
\[
 hH\subseteq H^2\subseteq H.
\]
Aplicația
\[
 \varphi_h:H\longrightarrow hH,\qquad
 \varphi_h(x)=hx
\]
este bijectivă. Deoarece $H$ este finită, $|hH|=|H|$, deci
\[
 hH=H^2=H.
\]
Cum $h\in H=hH$, există $x\in H$ cu $h=hx$, de unde
$x=1_G$. Astfel $1_G\in H$. Mai mult,
\[
 1_G\in hH\ \Longrightarrow\
  1_G=hx\ (x\in H)\ \Longrightarrow\
  x=h^{-1}\in H.
\]

\medskip
\noindent\underline{Teorema Cantor-Bernstein.}
Dacă $\alpha=|A|$, $\beta=|B|$, iar
\[
 \alpha\leq\beta,\qquad\beta\leq\alpha,
 \qquad\text{atunci}\qquad \alpha=\beta.
\]
; fără preambul sau pagină nouă.
% NOTE EDITORIALE:
% - Enunțul introductiv pare să compare monoizii multiplicativi
%   mN+1 și tN+1 cu U(Z_m), U(Z_t), probabil grupurile elementelor
%   inversabile ale inelelor de clase de resturi. Simbolul manuscris
%   seamănă cu U mare; această interpretare rămâne provizorie.
%   Pentru enunțul inițial nu există demonstrație pe această pagină.
% - În exercițiul 9, ENUNȚUL cere element neutru la stânga ex=x
%   și invers la DREAPTA xx'=e pentru fiecare x; combinația nu este
%   suficientă pentru un grup (există contraexemplu de semigrup drept-zero).
%   DEMONSTRAȚIA manuscrisă alege însă inverse la STÂNGA x'x=e,
%   x''x'=e și deduce xx'=e și xe=x. Păstrăm separat cele două
%   lecturi, fără corectarea tacită a enunțului.
% - La exercițiul 10 trebuie adăugată ipoteza H nevidă pentru
%   caracterizarea subgrupului; demonstrația presupune h în H.
% - Manuscrisul notează elementul neutru atât cu 1, cât și cu e.
% - Nu am reconstruit pasajele barate ori urmele de pe verso.

\subsection*{Seminar 4 - Pagina 4}

\noindent\underline{Arătați că:}
\[
\begin{gathered}
 (m\mathbb N+1,\cdot)\cong(t\mathbb N+1,\cdot),
 \qquad m,t\in\mathbb N \\[2pt]
 \Longleftrightarrow\\[-2pt]
 (\mathcal U(\mathbb Z_m),\cdot)
   \cong (\mathcal U(\mathbb Z_t),\cdot).
\end{gathered}
\]
\noindent (În manuscris, $\mathcal U(\mathbb Z_m)$ este indicat drept grup.)

\medskip
\noindent\underline{Exercițiul 9.}
Fie $(G,\cdot)$ semigrup. Atunci $G$ este grup dacă și numai dacă
\begin{enumerate}
 \item există $e\in G$ astfel încât $ex=x$, pentru orice $x\in G$;
 \item pentru orice $x\in G$ există $x'\in G$ cu
       $xx'=e$ (așa cum apare în enunț).
\end{enumerate}

\noindent\underline{Demonstrație (reciproca, după manuscris).}
Fie $x\in G$. Există $x'\in G$ astfel încât $x'x=e$ și
există $x''\in G$ astfel încât $x''x'=e$. Atunci
\[
\begin{aligned}
 xx'&=(ex)x'=(x''x')xx'
       =x''(x'x)x'=x''ex'=x''x'=e,\\
 xe&=x(x'x)=(xx')x=ex=x.
\end{aligned}
\]

\medskip
\noindent\underline{Exercițiul 10.}
Fie $(G,\cdot)$ grup, $H\subseteq G$, $H$ finită. Atunci
\[
 H\leq G\quad\Longleftrightarrow\quad
 H\text{ este parte stabilă la înmulțire}.
\]

\noindent\underline{Demonstrație.}
Implicația directă este evidentă. Reciproc, fie $h\in H$ și
\[
 hH=\{hx\mid x\in H\},\qquad
 H^2=\{xy\mid x,y\in H\}.
\]
Din stabilitatea lui $H$ rezultă
\[
 hH\subseteq H^2\subseteq H.
\]
Aplicația
\[
 \varphi_h:H\longrightarrow hH,\qquad
 \varphi_h(x)=hx
\]
este bijectivă. Deoarece $H$ este finită, $|hH|=|H|$, deci
\[
 hH=H^2=H.
\]
Cum $h\in H=hH$, există $x\in H$ cu $h=hx$, de unde
$x=1_G$. Astfel $1_G\in H$. Mai mult,
\[
 1_G\in hH\ \Longrightarrow\
  1_G=hx\ (x\in H)\ \Longrightarrow\
  x=h^{-1}\in H.
\]

\medskip
\noindent\underline{Teorema Cantor-Bernstein.}
Dacă $\alpha=|A|$, $\beta=|B|$, iar
\[
 \alpha\leq\beta,\qquad\beta\leq\alpha,
 \qquad\text{atunci}\qquad \alpha=\beta.
\]
; fără preambul sau pagină nouă.
% NOTE EDITORIALE:
% - Enunțul introductiv pare să compare monoizii multiplicativi
%   mN+1 și tN+1 cu U(Z_m), U(Z_t), probabil grupurile elementelor
%   inversabile ale inelelor de clase de resturi. Simbolul manuscris
%   seamănă cu U mare; această interpretare rămâne provizorie.
%   Pentru enunțul inițial nu există demonstrație pe această pagină.
% - În exercițiul 9, ENUNȚUL cere element neutru la stânga ex=x
%   și invers la DREAPTA xx'=e pentru fiecare x; combinația nu este
%   suficientă pentru un grup (există contraexemplu de semigrup drept-zero).
%   DEMONSTRAȚIA manuscrisă alege însă inverse la STÂNGA x'x=e,
%   x''x'=e și deduce xx'=e și xe=x. Păstrăm separat cele două
%   lecturi, fără corectarea tacită a enunțului.
% - La exercițiul 10 trebuie adăugată ipoteza H nevidă pentru
%   caracterizarea subgrupului; demonstrația presupune h în H.
% - Manuscrisul notează elementul neutru atât cu 1, cât și cu e.
% - Nu am reconstruit pasajele barate ori urmele de pe verso.

\subsection*{Seminar 4 - Pagina 4}

\noindent\underline{Arătați că:}
\[
\begin{gathered}
 (m\mathbb N+1,\cdot)\cong(t\mathbb N+1,\cdot),
 \qquad m,t\in\mathbb N \\[2pt]
 \Longleftrightarrow\\[-2pt]
 (\mathcal U(\mathbb Z_m),\cdot)
   \cong (\mathcal U(\mathbb Z_t),\cdot).
\end{gathered}
\]
\noindent (În manuscris, $\mathcal U(\mathbb Z_m)$ este indicat drept grup.)

\medskip
\noindent\underline{Exercițiul 9.}
Fie $(G,\cdot)$ semigrup. Atunci $G$ este grup dacă și numai dacă
\begin{enumerate}
 \item există $e\in G$ astfel încât $ex=x$, pentru orice $x\in G$;
 \item pentru orice $x\in G$ există $x'\in G$ cu
       $xx'=e$ (așa cum apare în enunț).
\end{enumerate}

\noindent\underline{Demonstrație (reciproca, după manuscris).}
Fie $x\in G$. Există $x'\in G$ astfel încât $x'x=e$ și
există $x''\in G$ astfel încât $x''x'=e$. Atunci
\[
\begin{aligned}
 xx'&=(ex)x'=(x''x')xx'
       =x''(x'x)x'=x''ex'=x''x'=e,\\
 xe&=x(x'x)=(xx')x=ex=x.
\end{aligned}
\]

\medskip
\noindent\underline{Exercițiul 10.}
Fie $(G,\cdot)$ grup, $H\subseteq G$, $H$ finită. Atunci
\[
 H\leq G\quad\Longleftrightarrow\quad
 H\text{ este parte stabilă la înmulțire}.
\]

\noindent\underline{Demonstrație.}
Implicația directă este evidentă. Reciproc, fie $h\in H$ și
\[
 hH=\{hx\mid x\in H\},\qquad
 H^2=\{xy\mid x,y\in H\}.
\]
Din stabilitatea lui $H$ rezultă
\[
 hH\subseteq H^2\subseteq H.
\]
Aplicația
\[
 \varphi_h:H\longrightarrow hH,\qquad
 \varphi_h(x)=hx
\]
este bijectivă. Deoarece $H$ este finită, $|hH|=|H|$, deci
\[
 hH=H^2=H.
\]
Cum $h\in H=hH$, există $x\in H$ cu $h=hx$, de unde
$x=1_G$. Astfel $1_G\in H$. Mai mult,
\[
 1_G\in hH\ \Longrightarrow\
  1_G=hx\ (x\in H)\ \Longrightarrow\
  x=h^{-1}\in H.
\]

\medskip
\noindent\underline{Teorema Cantor-Bernstein.}
Dacă $\alpha=|A|$, $\beta=|B|$, iar
\[
 \alpha\leq\beta,\qquad\beta\leq\alpha,
 \qquad\text{atunci}\qquad \alpha=\beta.
\]
; fără preambul sau pagină nouă.
% NOTE EDITORIALE:
% - Enunțul introductiv pare să compare monoizii multiplicativi
%   mN+1 și tN+1 cu U(Z_m), U(Z_t), probabil grupurile elementelor
%   inversabile ale inelelor de clase de resturi. Simbolul manuscris
%   seamănă cu U mare; această interpretare rămâne provizorie.
%   Pentru enunțul inițial nu există demonstrație pe această pagină.
% - În exercițiul 9, ENUNȚUL cere element neutru la stânga ex=x
%   și invers la DREAPTA xx'=e pentru fiecare x; combinația nu este
%   suficientă pentru un grup (există contraexemplu de semigrup drept-zero).
%   DEMONSTRAȚIA manuscrisă alege însă inverse la STÂNGA x'x=e,
%   x''x'=e și deduce xx'=e și xe=x. Păstrăm separat cele două
%   lecturi, fără corectarea tacită a enunțului.
% - La exercițiul 10 trebuie adăugată ipoteza H nevidă pentru
%   caracterizarea subgrupului; demonstrația presupune h în H.
% - Manuscrisul notează elementul neutru atât cu 1, cât și cu e.
% - Nu am reconstruit pasajele barate ori urmele de pe verso.

\subsection*{Seminar 4 - Pagina 4}

\noindent\underline{Arătați că:}
\[
\begin{gathered}
 (m\mathbb N+1,\cdot)\cong(t\mathbb N+1,\cdot),
 \qquad m,t\in\mathbb N \\[2pt]
 \Longleftrightarrow\\[-2pt]
 (\mathcal U(\mathbb Z_m),\cdot)
   \cong (\mathcal U(\mathbb Z_t),\cdot).
\end{gathered}
\]
\noindent (În manuscris, $\mathcal U(\mathbb Z_m)$ este indicat drept grup.)

\medskip
\noindent\underline{Exercițiul 9.}
Fie $(G,\cdot)$ semigrup. Atunci $G$ este grup dacă și numai dacă
\begin{enumerate}
 \item există $e\in G$ astfel încât $ex=x$, pentru orice $x\in G$;
 \item pentru orice $x\in G$ există $x'\in G$ cu
       $xx'=e$ (așa cum apare în enunț).
\end{enumerate}

\noindent\underline{Demonstrație (reciproca, după manuscris).}
Fie $x\in G$. Există $x'\in G$ astfel încât $x'x=e$ și
există $x''\in G$ astfel încât $x''x'=e$. Atunci
\[
\begin{aligned}
 xx'&=(ex)x'=(x''x')xx'
       =x''(x'x)x'=x''ex'=x''x'=e,\\
 xe&=x(x'x)=(xx')x=ex=x.
\end{aligned}
\]

\medskip
\noindent\underline{Exercițiul 10.}
Fie $(G,\cdot)$ grup, $H\subseteq G$, $H$ finită. Atunci
\[
 H\leq G\quad\Longleftrightarrow\quad
 H\text{ este parte stabilă la înmulțire}.
\]

\noindent\underline{Demonstrație.}
Implicația directă este evidentă. Reciproc, fie $h\in H$ și
\[
 hH=\{hx\mid x\in H\},\qquad
 H^2=\{xy\mid x,y\in H\}.
\]
Din stabilitatea lui $H$ rezultă
\[
 hH\subseteq H^2\subseteq H.
\]
Aplicația
\[
 \varphi_h:H\longrightarrow hH,\qquad
 \varphi_h(x)=hx
\]
este bijectivă. Deoarece $H$ este finită, $|hH|=|H|$, deci
\[
 hH=H^2=H.
\]
Cum $h\in H=hH$, există $x\in H$ cu $h=hx$, de unde
$x=1_G$. Astfel $1_G\in H$. Mai mult,
\[
 1_G\in hH\ \Longrightarrow\
  1_G=hx\ (x\in H)\ \Longrightarrow\
  x=h^{-1}\in H.
\]

\medskip
\noindent\underline{Teorema Cantor-Bernstein.}
Dacă $\alpha=|A|$, $\beta=|B|$, iar
\[
 \alpha\leq\beta,\qquad\beta\leq\alpha,
 \qquad\text{atunci}\qquad \alpha=\beta.
\]
